% method.tex
% ─────────────────────────────────────────────────────────────────────────────
% Supernode Abstraction for Circuit Tracing in Large Language Models:
% A Flow-Faithful Concept Grouping Framework
% ─────────────────────────────────────────────────────────────────────────────

\documentclass[11pt]{article}

\usepackage[margin=1.25in]{geometry}
\usepackage{amsmath, amsthm, amssymb}
\usepackage{mathtools}
\usepackage{bm}
\usepackage{bbm}
\usepackage{xcolor}
\usepackage{hyperref}
\usepackage{cleveref}
\usepackage{algorithm}
\usepackage{algpseudocode}
\usepackage{booktabs}
\usepackage{enumitem}
\usepackage{microtype}
\usepackage{natbib}

\hypersetup{colorlinks=true, linkcolor=blue!60!black, citecolor=green!50!black, urlcolor=blue}

% ── Theorem environments ──────────────────────────────────────────────────────
\newtheorem{theorem}{Theorem}[section]
\newtheorem{proposition}[theorem]{Proposition}
\newtheorem{lemma}[theorem]{Lemma}
\newtheorem{corollary}[theorem]{Corollary}
\newtheorem{definition}[theorem]{Definition}
\newtheorem{remark}[theorem]{Remark}
\newtheorem{assumption}[theorem]{Assumption}

% ── Notation macros ───────────────────────────────────────────────────────────
\newcommand{\cG}{\mathcal{G}}
\newcommand{\cF}{\mathcal{F}}
\newcommand{\cK}{\mathcal{K}}
\newcommand{\cL}{\mathcal{L}}
\newcommand{\cS}{\mathcal{S}}
\newcommand{\cV}{\mathcal{V}}
\newcommand{\cE}{\mathcal{E}}
\newcommand{\bA}{\mathbf{A}}
\newcommand{\bB}{\mathbf{B}}
\newcommand{\bC}{\mathbf{C}}
\newcommand{\bD}{\mathbf{D}}
\newcommand{\bL}{\mathbf{L}}
\newcommand{\bM}{\mathbf{M}}
\newcommand{\bP}{\mathbf{P}}
\newcommand{\bS}{\mathbf{S}}
\newcommand{\bW}{\mathbf{W}}
\newcommand{\bEps}{\mathbf{E}}
\newcommand{\btheta}{\bm{\theta}}
\newcommand{\bz}{\mathbf{z}}
\newcommand{\R}{\mathbb{R}}
\newcommand{\E}{\mathbb{E}}
\newcommand{\1}{\mathbbm{1}}
\newcommand{\eps}{\varepsilon}
\newcommand{\Dphi}{D(\varphi)}
\newcommand{\Rphi}{R(\varphi)}
\newcommand{\Rbal}{R_{\mathrm{bal}}(\varphi)}
\newcommand{\Rsup}{R_{\mathrm{sup}}(\varphi)}
\newcommand{\Fphi}{F(\varphi)}
\newcommand{\snAdj}{\mathbf{A}^{\mathrm{SN}}}
\newcommand{\snInf}{\mathbf{f}^{\mathrm{SN}}}
\newcommand{\snAct}{\mathbf{a}^{\mathrm{SN}}}
\newcommand{\norm}[1]{\left\lVert #1 \right\rVert}

% ─────────────────────────────────────────────────────────────────────────────

\title{\textbf{Flow-Faithful Supernode Abstraction for\\
Mechanistic Interpretation of LLM Circuits}\\[0.5em]
\large Method Description}

\author{}
\date{}

\begin{document}
\maketitle
\tableofcontents
\newpage

% ─────────────────────────────────────────────────────────────────────────────
\section{Introduction}
\label{sec:intro}
% ─────────────────────────────────────────────────────────────────────────────

Mechanistic interpretability of large language models (LLMs) aims to reverse-engineer the computational processes by which a model transforms an input prompt into an output prediction~\citep{elhage2021mathematical, wang2022interpretability}.
A central tool in this endeavor is \emph{circuit tracing}: extracting from the model's forward pass a sparse attribution graph $\cG = (\cV, \cE, w)$ whose nodes are individual features (typically transcoder or SAE latent dimensions) and whose directed edges carry attribution weights quantifying how strongly each upstream feature contributes to each downstream feature and ultimately to the logit of a target token.

Circuit graphs extracted from state-of-the-art LLMs routinely contain hundreds to thousands of nodes, rendering direct human interpretation intractable.
A natural remedy is \emph{abstraction}: grouping nodes that play similar computational roles into \emph{supernodes}, yielding a coarser graph $\cG^{\mathrm{SN}}$ that is small enough to inspect but ideally preserves the causal story encoded in $\cG$.

The challenge is that ``similar computational role'' is ambiguous.
Existing approaches either cluster by feature semantics (e.g., cosine similarity of decoder directions~\citep{cunningham2023sparse}) or by positional proximity in layer space, neither of which directly captures \emph{flow identity} --- the property that grouped nodes receive attribution from the same upstream sources and direct attribution to the same downstream targets.

\paragraph{Our contribution.}
We propose a principled supernode abstraction framework grounded in three ideas:
\begin{enumerate}[leftmargin=2em]
  \item \textbf{Structural equivalence clustering with mediation-aware similarity.}
        Nodes sharing similar weighted in- and out-neighbor profiles in $\cG$ likely represent the same abstract concept; we group them via spectral clustering on a cosine-similarity matrix derived from the attribution adjacency.
        Critically, we down-weight similarity between node pairs that have a mediating node strictly between their layers, preventing spectral clustering from grouping non-adjacent-layer nodes whose merger would introduce cycles in the supernode graph.

  \item \textbf{DAG-constrained layer alignment with containment cycle prevention.}
        We enforce that each supernode spans at most $\Delta$ consecutive layers and that no supernode's layer range contains another's, ensuring that cross-supernode edges represent genuine inter-concept information transfer rather than within-concept shortcuts.
        The latter condition---\emph{containment}---is a strictly stronger requirement than the interleaving check used in prior work and is necessary to guarantee acyclicity of $\cG^{\mathrm{SN}}$.

  \item \textbf{Flow faithfulness certification with suppression-aware residuals.}
        After clustering, we certify the quality of the abstraction with three grounded metrics---path concentration $\Dphi$, local flow residual $\Rphi$, and shortcut fraction $\sigma(\varphi)$---that together constitute a \emph{flow faithfulness score} $\Fphi \in [0,1]$.
        The local flow residual is decomposed into a \emph{balance residual} $\Rbal$, measuring genuine grouping quality for excitatory concept nodes, and a \emph{suppression ratio} $\Rsup$, measuring the inhibitory strength of suppressive concept nodes.
        Only $\Rbal$ enters the faithfulness score, preventing suppressive nodes from artifactually inflating the residual beyond its natural range.
        Correspondingly, the path attribution decomposition distributes flow only through positive-weight exits, so that suppressive edges do not distort the causal path enumeration.
\end{enumerate}

We provide theoretical analysis connecting each component to established results from community detection~\citep{abbe2018community} and directed graph clustering~\citep{rohe2016co}, and demonstrate on factual recall circuits extracted from Gemma-2-2B that the method produces interpretable, quantitatively certified concept-level causal graphs.

% ─────────────────────────────────────────────────────────────────────────────
\section{Problem Formulation}
\label{sec:problem}
% ─────────────────────────────────────────────────────────────────────────────

\subsection{Attribution Graphs}
\label{sec:attr-graph}

Let $M$ be an LLM and $x$ an input prompt.
Running a circuit tracing procedure (e.g., attribution patching~\citep{meng2022locating} or transcoders~\citep{templeton2024scaling}) yields a \emph{circuit attribution graph}:
\[
  \cG = (\cV,\, \cE,\, w),
\]
where:
\begin{itemize}[leftmargin=2em]
  \item $\cV = \cV_{\mathrm{emb}} \cup \cV_{\mathrm{mid}} \cup \cV_{\mathrm{logit}}$ partitions nodes into embedding nodes (input token representations), middle nodes (transcoder/SAE features at layers $1,\ldots,L-1$), and the target logit output node.
  \item $\cE \subseteq \cV \times \cV$ contains directed edges $(u,v)$ with $\mathrm{layer}(u) \le \mathrm{layer}(v)$, i.e., $\cG$ is a directed acyclic graph (DAG).
  \item $w: \cE \to \R$ assigns a signed attribution weight to each edge; $w(u,v) > 0$ indicates excitatory contribution and $w(u,v) < 0$ indicates suppressive contribution.
\end{itemize}

We write $N = |\cV|$, $N_{\mathrm{mid}} = |\cV_{\mathrm{mid}}|$, and $\bA \in \R^{N \times N}$ for the (sparse) adjacency matrix with $A_{ij} = w(i \to j)$.

Let $v^* \in \cV_{\mathrm{logit}}$ denote the target logit node.
The \emph{direct attribution} of node $i$ to the logit is $f_i = A_{i,v^*}$, i.e., the edge weight from $i$ directly to $v^*$.
The \emph{total logit-directed flow} is $F = \sum_{i \in \cV} f_i$.

\subsection{Supernode Abstraction}
\label{sec:sn-def}

\begin{definition}[Supernode Mapping]
  A \emph{supernode mapping} is a surjection $\varphi: \cV \to \{S_1, \ldots, S_K\}$ that assigns each node to one of $K$ supernodes, subject to:
  \begin{enumerate}[label=(\roman*)]
    \item $\varphi$ fixes embedding and logit nodes: each $v \in \cV_{\mathrm{emb}} \cup \cV_{\mathrm{logit}}$ forms its own singleton supernode.
    \item Each middle supernode $S_k$ has bounded layer span: $\max_{i \in S_k} \mathrm{layer}(i) - \min_{i \in S_k} \mathrm{layer}(i) \le \Delta$ for a fixed parameter $\Delta \ge 0$.
    \item The induced supernode graph $\cG^{\mathrm{SN}}$ is acyclic: supernode layer ranges neither interleave nor contain one another (see \Cref{def:dag-constraint}).
  \end{enumerate}
\end{definition}

\begin{definition}[DAG Constraint]
\label{def:dag-constraint}
  Two supernodes $S_k$ and $S_l$ with layer ranges $[\ell^-_k, \ell^+_k]$ and $[\ell^-_l, \ell^+_l]$ are \emph{compatible} if their ranges are strictly ordered or disjoint.
  They are \emph{incompatible} if their ranges \emph{interleave} ($\ell^-_k < \ell^-_l < \ell^+_k < \ell^+_l$ or vice versa) or if one \emph{contains} the other ($\ell^-_k < \ell^-_l$ and $\ell^+_l < \ell^+_k$, or vice versa).
\end{definition}

\begin{proposition}[Containment implies cycle]
\label{prop:containment-cycle}
  If $S_k$ contains $S_l$ (i.e., $\ell^-_k < \ell^-_l$ and $\ell^+_l < \ell^+_k$), then $A^{\mathrm{SN}}_{kl} > 0$ and $A^{\mathrm{SN}}_{lk} > 0$ whenever $S_k$ has members at layers both below $\ell^-_l$ and above $\ell^+_l$, creating a cycle in $\cG^{\mathrm{SN}}$.
\end{proposition}

\begin{proof}
  Let $i \in S_k$ with $\mathrm{layer}(i) < \ell^-_l$ and $i' \in S_k$ with $\mathrm{layer}(i') > \ell^+_l$, and let $j \in S_l$.
  Then $A_{ij}$ contributes to $A^{\mathrm{SN}}_{kl}$ (since $\mathrm{layer}(i) \le \mathrm{layer}(j)$, the forward mask passes), and $A_{ji'}$ contributes to $A^{\mathrm{SN}}_{lk}$ (since $\mathrm{layer}(j) \le \mathrm{layer}(i')$).
  Provided these edges are nonzero, both $A^{\mathrm{SN}}_{kl} > 0$ and $A^{\mathrm{SN}}_{lk} > 0$, i.e., $S_k \to S_l \to S_k$ is a directed cycle.
\end{proof}

Given $\varphi$, we define three \textbf{grounded supernode quantities} that partition the original graph quantities exactly:

\paragraph{Supernode activation.}
\[
  a^{\mathrm{SN}}(S_k) \;=\; \max_{i \in S_k} \mathrm{act}(i).
\]

\paragraph{Supernode logit flow.}
\[
  f^{\mathrm{SN}}(S_k) \;=\; \sum_{i \in S_k} A_{i,v^*}.
\]
By construction, $\sum_k f^{\mathrm{SN}}(S_k) = F$ (exact partition).

\paragraph{Supernode edge weight.}
\[
  A^{\mathrm{SN}}_{kl} \;=\; \sum_{\substack{i \in S_k,\, j \in S_l \\ \mathrm{layer}(i) \le \mathrm{layer}(j)}} A_{ij},
  \quad k \ne l.
\]
By construction, $\sum_{k \ne l} A^{\mathrm{SN}}_{kl}$ equals the sum of all cross-supernode forward edge weights in $\cG$ (exact partition).

\subsection{Desiderata}
\label{sec:desiderata}

A good supernode abstraction $\varphi$ should satisfy:
\begin{enumerate}[label=(D\arabic*), leftmargin=3em]
  \item \textbf{Conservation.} $\sum_k f^{\mathrm{SN}}(S_k) = F$ and edge weights are partitioned exactly. \emph{(Satisfied by construction.)}
  \item \textbf{Concept coherence.} Nodes within the same supernode have similar structural roles.
  \item \textbf{Flow faithfulness.} The dominant causal paths in $\cG$ project onto a small number of dominant paths in $\cG^{\mathrm{SN}}$, preserving the circuit's interpretive narrative.
  \item \textbf{DAG integrity.} The supernode graph is acyclic. This requires pairwise compatibility of all supernode layer ranges (\Cref{def:dag-constraint}): neither interleaving nor containment is permitted.
\end{enumerate}

(D1) and (D4) are enforced as hard constraints. (D2) is optimized by the clustering objective. (D3) is measured post-hoc and incorporated into the scoring objective.

% ─────────────────────────────────────────────────────────────────────────────
\section{Method Overview}
\label{sec:overview}
% ─────────────────────────────────────────────────────────────────────────────

The pipeline proceeds in four stages.

\begin{enumerate}[leftmargin=2em, label=\textbf{Stage \arabic*.}]
  \item \textbf{Similarity computation with mediation penalty.} Compute $\bS \in [0,1]^{N \times N}$ from the attribution adjacency and apply a mediation penalty to node pairs with an intermediate mediating node between their layers (\Cref{sec:similarity}).
  \item \textbf{Spectral clustering with DAG enforcement.} Cluster middle nodes via spectral clustering on $\bS$ to target $K$ groups, then enforce layer-span, interleaving, and containment constraints (\Cref{sec:clustering}).
  \item \textbf{Supernode graph construction.} Aggregate node-level quantities into the exact partition $(\snAdj, \snInf, \snAct)$ (\Cref{sec:sn-graph}).
  \item \textbf{Flow faithfulness assessment.} Compute $\Dphi$, $\Rbal$, $\Rsup$, $\sigma(\varphi)$, and $\Fphi$; use $\Fphi$ to guide $K$ selection via enhanced auto-$K$ search (\Cref{sec:flow-analysis}).
\end{enumerate}

% ─────────────────────────────────────────────────────────────────────────────
\section{Method Details}
\label{sec:details}
% ─────────────────────────────────────────────────────────────────────────────

\subsection{Node Similarity Matrix with Mediation Penalty}
\label{sec:similarity}

For each node $i$, define the weight matrix
\[
  \bW = \mathrm{diag}\!\left(\alpha \cdot \tilde{\mathbf{a}} + \beta \cdot \tilde{\mathbf{f}}\right), \quad \alpha, \beta \ge 0,\; \alpha + \beta = 1,
\]
with normalized activation $\tilde{\mathbf{a}}$ and influence $\tilde{\mathbf{f}}$.
Define out- and in-cosine similarity matrices:
\[
  \bS^{\mathrm{out}} = \mathrm{CosNorm}(\bA \bW \bA^\top), \qquad
  \bS^{\mathrm{in}}  = \mathrm{CosNorm}(\bA^\top \bW \bA),
\]
and the structural similarity $\bS^{\mathrm{struct}} = \frac{1}{2}\bS^{\mathrm{out}} + \frac{1}{2}\bS^{\mathrm{in}}$.

\paragraph{Mediation penalty.}
Grouping nodes $i$ (layer $\ell_i$) and $j$ (layer $\ell_j > \ell_i$) when a node $m$ exists with $\ell_i < \mathrm{layer}(m) < \ell_j$, $A_{im} > 0$, $A_{mj} > 0$ would create a cycle in $\cG^{\mathrm{SN}}$ (by \Cref{prop:containment-cycle}).
We penalize such pairs via:
\begin{equation}
  P_{ij} = \begin{cases}
    \gamma & \text{if } \exists\, m :\; \mathrm{layer}(i) < \mathrm{layer}(m) < \mathrm{layer}(j),\; A_{im} > 0,\; A_{mj} > 0, \\
    1      & \text{otherwise,}
  \end{cases}
  \label{eq:mediation-penalty}
\end{equation}
with penalty strength $\gamma \in (0, 1)$ (default $\gamma = 0.1$). The final similarity matrix is:
\begin{equation}
  S_{ij} = S^{\mathrm{struct}}_{ij} \cdot P_{ij}, \quad \bS \in [0,1]^{N \times N}.
  \label{eq:similarity}
\end{equation}

The penalty matrix $\bP$ is computed by iterating over distinct mediator layer values and accumulating binary mediator-existence matrices, giving $O(N^2 \cdot |\mathcal{L}|)$ time where $|\mathcal{L}|$ is the number of distinct layers.

\begin{remark}
  Setting $\gamma \to 0$ makes the penalty a hard block; $\gamma = 1$ disables it.
  In practice $\gamma = 0.1$ provides strong discouragement while remaining soft, with the DAG enforcement (\Cref{sec:clustering}) handling residual cases.
\end{remark}

\begin{remark}
  The use of both $\bS^{\mathrm{out}}$ and $\bS^{\mathrm{in}}$ corresponds to using both left and right singular vectors of $\bA$, as recommended for directed stochastic block models~\citep{rohe2016co}.
\end{remark}

\subsection{Spectral Clustering with DAG Enforcement}
\label{sec:clustering}

\paragraph{Eigengap heuristic.}
Let $\bL = \bD^{-1/2}(\bD - \bS_{\mathrm{mid}})\bD^{-1/2}$ be the normalized Laplacian of $\bS_{\mathrm{mid}}$.
We use $k^* = \arg\max_{k \ge 2}(\lambda_{k+1} - \lambda_k)$ as a heuristic~\citep{von2007tutorial}, defining search range $[k^* - 2, k^* + 2]$.

\paragraph{Spectral clustering.}
For target $K$, apply spectral clustering~\citep{ng2002spectral} to $\bS_{\mathrm{mid}}$ with $K$ clusters.
Outlier nodes (maximum similarity below global mean) are reassigned to the nearest cluster by mean layer distance.

\paragraph{DAG enforcement.}
Raw spectral clusters may still violate \Cref{def:dag-constraint} since the mediation penalty is soft.
\Cref{alg:dag} enforces DAG integrity by handling layer-span violations, interleaving, and containment.

\begin{algorithm}[t]
\caption{DAG Enforcement}
\label{alg:dag}
\begin{algorithmic}[1]
\Require Raw clusters $\{C_k\}$, layer map $\ell$, max span $\Delta$, max supernodes $K_{\max}$
\Ensure Supernode partition $\varphi$ satisfying DAG constraints
\State Initialize queue $Q \leftarrow \{C_k\}$, final list $\mathcal{F} \leftarrow \emptyset$
\While{$Q \ne \emptyset$} \Comment{Split oversized clusters}
  \State $C \leftarrow \mathrm{pop}(Q)$
  \If{$\max_C \ell - \min_C \ell \le \Delta$ \textbf{or} $|C| = 1$}
    \State $\mathcal{F} \leftarrow \mathcal{F} \cup \{C\}$
  \Else
    \State Split $C$ at midpoint layer; push both halves to $Q$
  \EndIf
\EndWhile
\Repeat \Comment{Resolve pairwise topology violations}
  \For{each pair $(F_a, F_b) \in \mathcal{F}^2$}
    \If{ranges interleave} \Comment{$\ell^-_a < \ell^-_b < \ell^+_a < \ell^+_b$ or symmetric}
      \State Split the wider cluster at the boundary of the narrower; restart
    \ElsIf{$\ell^-_a < \ell^-_b$ \textbf{and} $\ell^+_b < \ell^+_a$} \Comment{$F_a$ contains $F_b$}
      \State Split $F_a$ at $\ell^-_b$ into $F_a^{\mathrm{lo}} = \{i \in F_a : \ell(i) < \ell^-_b\}$ and $F_a^{\mathrm{hi}} = \{i \in F_a : \ell(i) \ge \ell^-_b\}$; restart
    \ElsIf{$\ell^-_b < \ell^-_a$ \textbf{and} $\ell^+_a < \ell^+_b$} \Comment{$F_b$ contains $F_a$}
      \State Split $F_b$ at $\ell^-_a$ analogously; restart
    \EndIf
  \EndFor
\Until{no violations remain}
\If{$|\mathcal{F}| > K_{\max} - n_{\mathrm{fixed}}$}
  \State Greedily merge closest-centroid pairs until budget met
\EndIf
\State \Return $\varphi$ with singleton supernodes for $\cV_{\mathrm{emb}} \cup \cV_{\mathrm{logit}}$
\end{algorithmic}
\end{algorithm}

\begin{proposition}[Correctness of DAG Enforcement]
\label{prop:dag-correctness}
  \Cref{alg:dag} terminates and produces a partition where all supernode pairs are compatible (\Cref{def:dag-constraint}).
\end{proposition}

\begin{proof}
  \emph{Termination.} Each split strictly reduces the layer span of at least one cluster and increases $|\mathcal{F}|$ by one, bounded above by $N_{\mathrm{mid}}$.
  \emph{Correctness.} After the repeat loop, no pair interleaves (by the interleaving branch) and no pair has one containing the other (by the containment branches), leaving only strict ordering or disjointness.
\end{proof}

\begin{remark}
  The containment split at $\ell^-_b$ ensures $F_a^{\mathrm{lo}}$ lies strictly below $F_b$.
  If $F_a^{\mathrm{hi}}$ still contains $F_b$ (i.e., $\ell^+_b < \ell^+_a$), a subsequent iteration splits $F_a^{\mathrm{hi}}$ at $\ell^+_b$, handling arbitrary nesting depth.
\end{remark}

\subsection{Supernode Graph Construction}
\label{sec:sn-graph}

Given $\varphi$, construct $\cG^{\mathrm{SN}}$ with quantities as defined in \Cref{sec:sn-def}.
The conservation identities
\begin{equation}
  \sum_{k=1}^K f^{\mathrm{SN}}(S_k) = F, \qquad
  \sum_{k \ne l} A^{\mathrm{SN}}_{kl} = \sum_{\substack{(i,j) \in \cE \\ \varphi(i) \ne \varphi(j)}} A_{ij}
  \label{eq:conservation}
\end{equation}
hold to floating-point precision by construction and are verified empirically after each grouping.

\subsection{Flow Faithfulness Analysis}
\label{sec:flow-analysis}

Conservation (\Cref{eq:conservation}) is necessary but not sufficient for a faithful abstraction.
We introduce three complementary metrics.

\subsubsection{Path Attribution Decomposition and $\Dphi$}
\label{sec:path-decomp}

A \emph{supernode path} is a sequence $\pi = (S_{e}, S_{a_1}, \ldots, S_{a_m}, S_{\ell})$ from an embedding supernode to the logit supernode.
We propagate flow forward through the DAG $\cG^{\mathrm{SN}}$ using proportional distribution over \emph{positive} exits only:
\begin{equation}
  w(\pi_{1:t} \to S_b) \;=\; w(\pi_{1:t}) \cdot \frac{A^{\mathrm{SN}}_{\pi_t, S_b}}{\displaystyle\sum_{S'} \max(0,\, A^{\mathrm{SN}}_{\pi_t, S'}) + \max(0,\, f^{\mathrm{SN}}(\pi_t))},
  \label{eq:flow-prop}
\end{equation}
where the denominator sums only positive exit weights (positive edges out plus positive direct logit flow).
Suppressive exits ($A^{\mathrm{SN}} < 0$ or $f^{\mathrm{SN}} < 0$) are excluded from the denominator: they represent inhibition rather than signal routing and should not distort the proportional flow distribution.
Since $\cG^{\mathrm{SN}}$ is a DAG (guaranteed by \Cref{alg:dag}), this forward pass terminates in $O(K^2)$ steps.

\begin{definition}[Flow Distortion]
  \begin{equation}
    D(\varphi) = 1 - \frac{\sum_{\pi \in \mathrm{top}_k \Pi} w(\pi)}{\sum_{\pi \in \Pi} w(\pi)},
    \label{eq:D-phi}
  \end{equation}
  where $\Pi$ is the set of all supernode paths with positive flow and $\mathrm{top}_k \Pi$ denotes the $k$ highest-weight paths.
\end{definition}

\subsubsection{Local Flow Residuals and $\Rbal$, $\Rsup$}
\label{sec:residuals}

For each middle supernode $S_k$, define signed net in- and out-flow:
\begin{align}
  \mathrm{in}(S_k) &= \sum_{l \ne k} A^{\mathrm{SN}}_{lk}, \label{eq:inflow} \\
  \mathrm{out}(S_k) &= \sum_{l \ne k} A^{\mathrm{SN}}_{kl} + f^{\mathrm{SN}}(S_k), \label{eq:outflow}
\end{align}
and the relative residual $r_k = |\mathrm{in}(S_k) - \mathrm{out}(S_k)| / (|\mathrm{in}(S_k)| + \eps)$.

\paragraph{Suppressive concept nodes.}
A supernode $S_k$ is \emph{suppressive} if $f^{\mathrm{SN}}(S_k) < 0$.
For such nodes, a large $r_k$ can arise because the inhibitory exit is mechanistically meaningful---the signed flow is a circuit property, not a grouping artifact.
Concretely, if $\mathrm{in}(S_k) > 0$ and $f^{\mathrm{SN}}(S_k) \ll 0$, then $\mathrm{out}(S_k) < 0$ and $r_k > 1$, inflating any global average beyond its natural $[0,1]$ range without indicating a grouping problem.

We therefore partition middle supernodes into:
\[
  \cV^{\mathrm{bal}}_{\mathrm{mid}} = \{S_k : f^{\mathrm{SN}}(S_k) \ge 0\}, \quad
  \cV^{\mathrm{sup}}_{\mathrm{mid}} = \{S_k : f^{\mathrm{SN}}(S_k) < 0\},
\]
and define two independent metrics:

\begin{definition}[Balance Residual]
  \begin{equation}
    \Rbal = \frac{1}{|\cV^{\mathrm{bal}}_{\mathrm{mid}}|} \sum_{S_k \in \cV^{\mathrm{bal}}_{\mathrm{mid}}} r_k.
    \label{eq:R-bal}
  \end{equation}
  $\Rbal \approx 0$ certifies that flow entering each excitatory concept roughly equals the flow leaving it.
\end{definition}

\begin{definition}[Suppression Ratio]
  \begin{equation}
    \Rsup = \frac{1}{|\cV^{\mathrm{sup}}_{\mathrm{mid}}|} \sum_{S_k \in \cV^{\mathrm{sup}}_{\mathrm{mid}}} \frac{|f^{\mathrm{SN}}(S_k)|}{|\mathrm{in}(S_k)| + \eps}.
    \label{eq:R-sup}
  \end{equation}
  $\Rsup$ characterizes the inhibitory strength of suppressive concepts and is reported as a diagnostic; it does not enter $\Fphi$.
\end{definition}

\begin{remark}
  A high $\Rbal$ indicates a fixable grouping problem (flow enters a concept and does not leave).
  A high $\Rsup$ indicates a circuit-level property (the model actively suppresses the target prediction).
  Conflating the two, as a single $R(\varphi)$ over all middle supernodes would do, masks both signals.
\end{remark}

\subsubsection{Shortcut Analysis and $\sigma(\varphi)$}
\label{sec:shortcuts}

\begin{definition}[Shortcut Ratio]
  For each edge $(S_a, S_c)$ with $A^{\mathrm{SN}}_{ac} > 0$:
  \begin{equation}
    \rho(S_a, S_c) = \frac{A^{\mathrm{SN}}_{ac}}{A^{\mathrm{SN}}_{ac} + \displaystyle\max_{S_b \ne S_a, S_c} \min\!\left(A^{\mathrm{SN}}_{ab},\, A^{\mathrm{SN}}_{bc}\right)}.
    \label{eq:shortcut-ratio}
  \end{equation}
\end{definition}

An edge is a shortcut if $\rho < 0.5$.

\begin{definition}[Global Shortcut Fraction]
  \begin{equation}
    \sigma(\varphi) = \frac{\sum_{(a,c):\, \rho < 0.5} A^{\mathrm{SN}}_{ac}}{\sum_{(a,c)} A^{\mathrm{SN}}_{ac}}.
    \label{eq:sigma-phi}
  \end{equation}
\end{definition}

\subsubsection{Combined Flow Faithfulness Score}
\label{sec:combined}

\begin{definition}[Flow Faithfulness Score]
  \begin{equation}
    F(\varphi) = w_p \cdot (1 - D(\varphi)) + w_r \cdot (1 - \Rbal) + w_s \cdot (1 - \sigma(\varphi)),
    \label{eq:F-phi}
  \end{equation}
  with default weights $w_p = 0.4$, $w_r = 0.3$, $w_s = 0.3$.
\end{definition}

$\Fphi \in [0,1]$ certifies abstraction quality. $\Rsup$ is reported alongside but does not affect the score, since suppressive behavior is a mechanistic circuit property rather than an abstraction artifact.

\subsubsection{Enhanced Auto-$K$ Selection}
\label{sec:auto-k}

\begin{equation}
  \hat{K} = \arg\max_{K \in [k^*-2,\, k^*+2]} \cL(K),
  \label{eq:auto-k}
\end{equation}
\begin{equation}
  \cL(K) = w_1 \cdot \bar{s}_{\mathrm{intra}}(K) + w_2 \cdot \mathrm{DAG}(K) + w_3 \cdot H_{\mathrm{attr}}(K) + w_4 \cdot \mathrm{size}(K) + w_5 \cdot F(\varphi_K),
  \label{eq:score}
\end{equation}
where $\bar{s}_{\mathrm{intra}}$ is weighted mean intra-supernode similarity, $\mathrm{DAG}(K) = 1 - n_{\mathrm{warn}}/\binom{K}{2}$, $H_{\mathrm{attr}}$ is normalized entropy of $f^{\mathrm{SN}}$, and $\mathrm{size}(K) = 1 - |K - \sqrt{N_{\mathrm{mid}}}|/N_{\mathrm{mid}}$.

% ─────────────────────────────────────────────────────────────────────────────
\section{Theoretical Analysis}
\label{sec:theory}
% ─────────────────────────────────────────────────────────────────────────────

\subsection{Generative Model: Layered Concept Block Model}
\label{sec:model}

\begin{assumption}[Layered Degree-Corrected Block Model (LDCBM)]
\label{ass:block}
  There exist $K$ latent concept groups with assignments $z_i \in [K]$.
  \begin{equation}
    \E[A_{ij}] = \theta_i \theta_j B_{z_i, z_j} \cdot \1[\mathrm{layer}(i) < \mathrm{layer}(j)],
    \label{eq:ldcbm}
  \end{equation}
  where $\theta_i > 0$ and $\bB \in \R^{K \times K}_{\ge 0}$.
  Observed weights: $A_{ij} = \E[A_{ij}] + \eps_{ij}$, $\mathrm{Var}(\eps_{ij}) \le \sigma^2$.
\end{assumption}

\begin{assumption}[Identifiability]
\label{ass:ident}
  $B_{k\cdot} \ne B_{l\cdot}$ for all $k \ne l$.
\end{assumption}

\begin{assumption}[Degree regularity]
\label{ass:degree}
  $\bar{d} = \omega(\log N_{\mathrm{mid}})$.
\end{assumption}

\subsection{Cycle Prevention under the LDCBM}
\label{sec:cycle-theory}

\begin{proposition}[Mediation penalty eliminates cross-block containment]
\label{prop:mediation-prevents-cycle}
  Under \Cref{ass:block}, suppose $z_i \ne z_j$, $\mathrm{layer}(i) < \mathrm{layer}(j)$, and there exists $m$ with $z_m \notin \{z_i, z_j\}$, $\mathrm{layer}(i) < \mathrm{layer}(m) < \mathrm{layer}(j)$, $\E[A_{im}] > 0$, $\E[A_{mj}] > 0$.
  Then $P_{ij} = \gamma$ (\Cref{eq:mediation-penalty}).
  As $\gamma \to 0$, the probability that spectral clustering assigns $i$ and $j$ to the same cluster converges to zero.
\end{proposition}

\begin{proof}
  Under the LDCBM, $\E[S^{\mathrm{struct}}_{ii'}] = C_{z_i z_i} > C_{z_i z_j}$ for same-block nodes (by \Cref{ass:ident}).
  The penalized similarity $\gamma \cdot S^{\mathrm{struct}}_{ij}$ falls below within-block similarities when $\gamma < C_{z_i z_j} / C_{z_i z_i}$, causing spectral clustering to assign $i$ and $j$ to different clusters with high probability, preventing the containment cycle identified in \Cref{prop:containment-cycle}.
\end{proof}

\subsection{Concept Recovery Guarantee}
\label{sec:recovery}

\begin{theorem}[Concept Recovery]
\label{thm:recovery}
  Under \Cref{ass:block,ass:ident,ass:degree}, if $\mathrm{SNR} \ge C \sqrt{\log N_{\mathrm{mid}} / \bar{d}}$, then spectral clustering on $\bS_{\mathrm{mid}}$ recovers the true partition with probability $1 - o(1)$, up to $o(N_{\mathrm{mid}})$ errors.
\end{theorem}

\begin{proof}
  \emph{Step 1.} Under the LDCBM, $\E[\bS^{\mathrm{struct}}_{\mathrm{mid}}]_{ij} = C_{z_i z_j}$ (block-constant), since $\theta_i\theta_j$ factors cancel in cosine normalization. The mediation penalty preserves block structure: same-block nodes share layers (no mediating paths) so $P_{ij} = 1$ within blocks; cross-block penalized entries remain block-constant with enlarged between-block gap.

  \emph{Step 2.} By matrix Bernstein~\citep{tropp2012user}: $\|\bS_{\mathrm{mid}} - \E[\bS_{\mathrm{mid}}]\|_2 \le c\sqrt{\log N_{\mathrm{mid}}/\bar{d}}$ w.p. $1 - N_{\mathrm{mid}}^{-2}$.

  \emph{Step 3.} The spectral gap $\lambda_K - \lambda_{K+1} \ge c' \cdot \mathrm{SNR}$ (Davis-Kahan~\citep{davis1970rotation}).

  \emph{Step 4.} By \citet{lei2015consistency}, when $\mathrm{SNR} \ge C\sqrt{\log N / \bar{d}}$, at most $o(N_{\mathrm{mid}})$ nodes are misclassified.

  \emph{Step 5.} The layer-span and containment constraints of \Cref{alg:dag} eliminate spurious cross-layer groupings without affecting correctly classified same-block assignments (which have compatible layer ranges under \Cref{ass:block}).
\end{proof}

\subsection{Flow Faithfulness Under Correct Recovery}
\label{sec:flow-theory}

\begin{theorem}[Flow Distortion Bound]
\label{thm:flow-distortion}
  Under correct recovery:
  \begin{equation}
    D(\varphi) \;\le\; \frac{W_{\mathrm{sc}}}{W_{\mathrm{total}}},
  \end{equation}
  where $W_{\mathrm{total}} = \sum_{(a,c)} \max(0, A^{\mathrm{SN}}_{ac})$ (total positive edge weight) and $W_{\mathrm{sc}}$ is the weight of shortcut edges.
\end{theorem}

\begin{proof}
  Path propagation (\Cref{eq:flow-prop}) distributes only positive exit flow, so $\sum_\pi w(\pi) = W_{\mathrm{total}}$.
  Shortcut edges disperse flow across non-dominant paths, preventing top-$k$ concentration.
  The fraction misrouted by shortcuts is bounded by $W_{\mathrm{sc}}/W_{\mathrm{total}}$.
\end{proof}

\begin{theorem}[Local Conservation Bound]
\label{thm:local-residual}
  Under \Cref{ass:block} with correct recovery, for $S_k \in \cV^{\mathrm{bal}}_{\mathrm{mid}}$:
  \begin{equation}
    r_k \;\le\; \frac{\delta_k}{\bar{d}^{\mathrm{inter}}_k + \eps},
  \end{equation}
  where $\delta_k = \max_{i,j \in S_k}|\sum_m A_{im}/\theta_i - \sum_m A_{jm}/\theta_j|$ is within-block flow asymmetry and $\bar{d}^{\mathrm{inter}}_k$ is mean inter-block out-degree. Consequently:
  \begin{equation}
    \Rbal \;\le\; \frac{\max_{k \in \cV^{\mathrm{bal}}_{\mathrm{mid}}} \delta_k}{\min_{k \in \cV^{\mathrm{bal}}_{\mathrm{mid}}} \bar{d}^{\mathrm{inter}}_k + \eps}.
    \label{eq:R-phi-bound}
  \end{equation}
  The bound applies only to $\cV^{\mathrm{bal}}_{\mathrm{mid}}$; suppressive nodes are excluded because their signed exit is a circuit property, not a grouping artifact.
\end{theorem}

\begin{proof}
  Under the LDCBM, homogeneous blocks ($\delta_k = 0$) satisfy $\mathrm{in}(S_k) = \mathrm{out}(S_k)$ exactly, giving $r_k = 0$.
  In the heterogeneous case, $|\mathrm{in}(S_k) - \mathrm{out}(S_k)| \le |S_k| \cdot \delta_k \cdot \bar\theta_k$ and $|\mathrm{in}(S_k)| \ge |S_k| \cdot \bar{d}^{\mathrm{inter}}_k$, yielding $r_k \le \delta_k / \bar{d}^{\mathrm{inter}}_k$. Averaging over $\cV^{\mathrm{bal}}_{\mathrm{mid}}$ gives \Cref{eq:R-phi-bound}.
\end{proof}

\begin{proposition}[Path Concentration]
\label{prop:path-concentration}
  If $\cG^{\mathrm{SN}}$ has positive out-degree $\le d$ and DAG depth $\le L$, then:
  \begin{equation}
    D(\varphi) \;\le\; \frac{d^L - k}{d^L}.
  \end{equation}
  Suppressive edges do not increase path count since they are excluded from flow propagation (\Cref{eq:flow-prop}).
\end{proposition}

\subsection{Combined Guarantee}
\label{sec:combined-guarantee}

\begin{theorem}[Flow Faithfulness Certificate]
\label{thm:main}
  Under \Cref{ass:block,ass:ident,ass:degree} with sufficient SNR, the pipeline produces $\varphi$ satisfying with probability $1 - o(1)$:
  \begin{enumerate}[label=(\alph*)]
    \item \emph{(Recovery)} Concept partition recovered up to $o(N)$ errors.
    \item \emph{(Cycle-free)} $\cG^{\mathrm{SN}}$ is a DAG: mediation penalty prevents containment groupings at clustering (\Cref{prop:mediation-prevents-cycle}); \Cref{alg:dag} eliminates residual violations (\Cref{prop:dag-correctness}).
    \item \emph{(Flow distortion)} $D(\varphi) \le W_{\mathrm{sc}} / W_{\mathrm{total}}$ (positive edges only).
    \item \emph{(Local conservation)} $\Rbal \le \max_k \delta_k / \min_k \bar{d}^{\mathrm{inter}}_k$ over $\cV^{\mathrm{bal}}_{\mathrm{mid}}$; suppressive nodes characterized separately by $\Rsup$.
    \item \emph{(Path concentration)} $D(\varphi) \le (d^L - k)/d^L$ for positive out-degree $\le d$, depth $\le L$.
  \end{enumerate}
  $D(\varphi)$, $\Rbal$, and $\Rsup$ are all computable from $\cG^{\mathrm{SN}}$ without ground-truth knowledge, providing model-free validity certificates.
\end{theorem}

\begin{proof}
  (a) \Cref{thm:recovery}. (b) \Cref{prop:mediation-prevents-cycle,prop:dag-correctness}. (c) \Cref{thm:flow-distortion}. (d) \Cref{thm:local-residual}. (e) \Cref{prop:path-concentration}.
\end{proof}

% ─────────────────────────────────────────────────────────────────────────────
\section{Discussion and Limitations}
\label{sec:discussion}
% ─────────────────────────────────────────────────────────────────────────────

\paragraph{Suppressive concept nodes.}
The framework distinguishes suppressive supernodes ($f^{\mathrm{SN}}(S_k) < 0$) from excitatory ones.
Suppressive nodes are characterized by $\Rsup$ and excluded from both $\Rbal$ and the path flow propagation (\Cref{eq:flow-prop}).
This separation is essential: a large $|\mathrm{out}|/|\mathrm{in}|$ ratio for a suppressive node reflects the model actively inhibiting competing predictions---a mechanistically meaningful signal that would be obscured by unsigned aggregation or conflated with grouping artifacts in a single $R(\varphi)$ metric.

\paragraph{Containment cycles and mediation penalty.}
A subtle failure mode of naive clustering is grouping nodes $A$ (layer $l_1$) and $C$ (layer $l_3 > l_1$) when a mediating node $B$ (layer $l_2 \in (l_1, l_3)$) carries flow $A \to B \to C$.
This creates cycle $\{A,C\} \to B \to \{A,C\}$ (\Cref{prop:containment-cycle}), invalidating the DAG property and corrupting path propagation.
The mediation penalty (\Cref{eq:mediation-penalty}) addresses this at the clustering stage; \Cref{alg:dag} provides the hard guarantee via post-hoc containment splitting.

\paragraph{Sink supernodes and flow distortion.}
A supernode with large positive in-flow but near-zero out-flow and zero direct logit connection is a \emph{flow sink}, flagged by $r_k \approx 1$.
This indicates either nodes that fire but are causally irrelevant for this prediction, or over-aggressive circuit pruning.
Such nodes require manual inspection.

\paragraph{Block model approximation.}
The LDCBM is an idealization; real circuits may exhibit polysemanticity, hierarchical concept structure, or non-stationarity across layers.
The metrics $D(\varphi)$, $\Rbal$, and $\Rsup$ serve as model-free certificates regardless of whether the generative model holds exactly.

\paragraph{Computational complexity.}
Similarity with mediation penalty: $O(N^2 \cdot |\mathcal{L}|)$.
Spectral clustering: $O(N^2 K)$.
DAG enforcement: $O(N_{\mathrm{mid}}^2 L)$ worst case.
Path decomposition: $O(K^2 L)$.
All stages are tractable for $N \le 10^3$.

% ─────────────────────────────────────────────────────────────────────────────
\bibliographystyle{plainnat}
\begin{thebibliography}{99}

\bibitem[Abbe(2018)]{abbe2018community}
Abbe, E. (2018).
\newblock Community detection and stochastic block models: Recent developments.
\newblock \emph{Journal of Machine Learning Research}, 18(177):1--86.

\bibitem[Davis and Kahan(1970)]{davis1970rotation}
Davis, C. and Kahan, W.~M. (1970).
\newblock The rotation of eigenvectors by a perturbation. III.
\newblock \emph{SIAM Journal on Numerical Analysis}, 7(1):1--46.

\bibitem[Elhage et al.(2021)]{elhage2021mathematical}
Elhage, N., Nanda, N., Olsson, C., et al. (2021).
\newblock A mathematical framework for transformer circuits.
\newblock \emph{Transformer Circuits Thread}.

\bibitem[Lei and Rinaldo(2015)]{lei2015consistency}
Lei, J. and Rinaldo, A. (2015).
\newblock Consistency of spectral clustering in stochastic block models.
\newblock \emph{Annals of Statistics}, 43(1):215--237.

\bibitem[Meng et al.(2022)]{meng2022locating}
Meng, K., Bau, D., Andonian, A., and Belinkov, Y. (2022).
\newblock Locating and editing factual associations in {GPT}.
\newblock \emph{Advances in Neural Information Processing Systems}, 35.

\bibitem[Ng et al.(2002)]{ng2002spectral}
Ng, A., Jordan, M., and Weiss, Y. (2002).
\newblock On spectral clustering: Analysis and an algorithm.
\newblock \emph{Advances in Neural Information Processing Systems}, 14.

\bibitem[Qin and Rohe(2013)]{qin2013regularized}
Qin, T. and Rohe, K. (2013).
\newblock Regularized spectral clustering under the degree-corrected stochastic block model.
\newblock \emph{Advances in Neural Information Processing Systems}, 26.

\bibitem[Rohe et al.(2016)]{rohe2016co}
Rohe, K., Qin, T., and Yu, B. (2016).
\newblock Co-clustering directed graphs to discover asymmetries and directional communities.
\newblock \emph{Proceedings of the National Academy of Sciences}, 113(45):12679--12684.

\bibitem[Templeton et al.(2024)]{templeton2024scaling}
Templeton, A., Conerly, T., Marcus, J., et al. (2024).
\newblock Scaling monosemanticity: Extracting interpretable features from {Claude} {Sonnet}.
\newblock \emph{Transformer Circuits Thread}.

\bibitem[Tropp(2012)]{tropp2012user}
Tropp, J.~A. (2012).
\newblock User-friendly tail bounds for sums of random matrices.
\newblock \emph{Foundations of Computational Mathematics}, 12(4):389--434.

\bibitem[von Luxburg(2007)]{von2007tutorial}
von Luxburg, U. (2007).
\newblock A tutorial on spectral clustering.
\newblock \emph{Statistics and Computing}, 17(4):395--416.

\bibitem[Wang et al.(2022)]{wang2022interpretability}
Wang, K., Variengien, A., Conmy, A., Shlegeris, B., and Steinhardt, J. (2022).
\newblock Interpretability in the wild: A circuit for indirect object identification in {GPT-2} small.
\newblock \emph{arXiv preprint arXiv:2211.00593}.

\end{thebibliography}

\end{document}